\documentclass[11pt,a4paper]{article}
\usepackage[T1]{fontenc}
\usepackage{lmodern,amsmath,amssymb,booktabs,geometry,enumitem,fancyhdr,microtype,hyperref}
\geometry{margin=22mm}
\setlength{\parindent}{0pt}
\setlength{\parskip}{6pt}
\pagestyle{fancy}
\fancyhf{}
\fancyhead[L]{Machine Learning Club}
\fancyhead[R]{Advanced 01}
\fancyfoot[C]{\thepage}
\newcommand{\problem}[1]{\subsection*{Problem #1}}
\newcommand{\worked}[1]{\subsection*{Worked Example #1}}
\begin{document}
{\small SEVEN HILLS INTERNATIONAL SCHOOL \hfill 2026--2027}

{\Large\bfseries What can examples actually prove?}\par
Progressive algorithmic reasoning about generalisation, alternative rules and missing observations.

\section*{Preliminaries}
Let an input be a pair of Boolean features $(x,y)\in\{0,1\}^2$. A classifier returns a label from $\{0,1\}$. A classifier is \emph{consistent} with a training sample when it produces the supplied labels on all observed examples. There are no assumptions about the classifier's behaviour on unseen examples. Ordinary truth tables and pseudocode are sufficient.

\worked{A --- Two perfect models}
Our training set contains $(0,0)\mapsto 0$ and $(1,1)\mapsto 1$. Consider:
\begin{verbatim}
def model_a(x, y): return x
def model_b(x, y): return y
\end{verbatim}
\begin{center}
\begin{tabular}{ccccc}\toprule
$x$ & $y$ & model\_a & model\_b & observed?\\\midrule
0&0&0&0&yes\\
0&1&0&1&no\\
1&0&1&0&no\\
1&1&1&1&yes\\\bottomrule
\end{tabular}
\end{center}
\textbf{Solution.} Both models have perfect training accuracy but disagree on the two unobserved inputs. Thus perfect agreement with a training sample need not determine the true output on an unseen case.

\problem{1 --- How many models are possible?}
A binary classifier may choose its label independently on each of the four inputs.
\begin{enumerate}[label=(\alph*)]
\item How many classifiers are consistent with the sample above?
\item Justify your count without listing every classifier.
\item How many remain after one additional previously unobserved input is labelled?
\end{enumerate}
\textbf{Selected solution.} Two distinct inputs are still unobserved. Each may be labelled 0 or 1 independently, so $2^2=4$ consistent functions remain. After one new labelled input, $2$ remain.

\problem{2 --- A distinguishing example}
Suppose exactly two candidate target models remain: \texttt{model\_a} and \texttt{model\_b}.
\begin{enumerate}[label=(\alph*)]
\item Find one previously unobserved input on which they disagree.
\item Explain how its true label identifies one of the two candidates.
\item Explain why this would not establish correctness if the set of possible hypotheses were larger.
\end{enumerate}
\textbf{Hint.} Find a row with different outputs in the truth table.

\problem{3 --- An irrelevant feature}
Introduce a third Boolean input $z$, with the target output still determined by $x$ and $y$.
\begin{enumerate}[label=(\alph*)]
\item Give a pair of inputs that differ only in $z$ and must have equal outputs.
\item Explain how tests of $z$ can be removed from a tree that computes a function independent of $z$ on the entire domain.
\item Why can incomplete training data make $z$ appear useful even when it is irrelevant?
\end{enumerate}
\textbf{Hint.} Fix $z=0$ consistently in an equivalent full-domain representation.

\problem{4 --- A generalisation counterexample}
Suppose all inputs except $u$ have been observed and the learner must predict the label at $u$.
\begin{enumerate}[label=(\alph*)]
\item Construct two target functions identical on every observed input but opposite at $u$.
\item Explain why no deterministic learner can guarantee a correct label at $u$ without further assumptions.
\item Describe one useful additional assumption in an actual ML problem.
\end{enumerate}
\textbf{Hint.} Start from any target function and change only its label at $u$.

\section*{Observation}
Neither perfect training accuracy nor a large number of training examples logically establishes a unique label for an unobserved input without assumptions restricting the model or the data-generating process. Optional computational assistance: enumerate the four Boolean inputs in Python. No special software is needed.

\section*{References}
Experience AI, \emph{Bias in, bias out}; IOAI classical ML foundations.
\end{document}
